% ============================================================
% LEY DE COULOMB
% ============================================================


% ============================================================
% FUERZA ELÉCTRICA
% ============================================================

\section{Fuerza eléctrica}

Sabemos experimentalmente que dos cuerpos cargados pueden atraerse o
repelerse.

Hasta ahora esta interacción se ha descrito de manera cualitativa:

\[
\text{cargas del mismo signo}
\Longrightarrow
\text{repulsión},
\]

mientras que

\[
\text{cargas de signos opuestos}
\Longrightarrow
\text{atracción}.
\]

Sin embargo, para describir físicamente una interacción no basta con saber
si existe atracción o repulsión.

También debemos poder determinar:

\begin{itemize}
    \item cuánto vale la fuerza;
    \item en qué dirección actúa;
    \item cuál es su sentido;
    \item cómo cambia cuando modificamos las cargas;
    \item cómo cambia cuando modificamos la distancia entre ellas.
\end{itemize}

La relación que permite responder estas preguntas para cargas puntuales es
la \textbf{ley de Coulomb}.


% ============================================================
% CONCEPTO DE FUERZA ELÉCTRICA
% ============================================================

\subsection{La fuerza eléctrica como interacción}

\begin{definicion}[Fuerza eléctrica]

La \textbf{fuerza eléctrica} es la interacción de atracción o repulsión que
puede producirse entre cuerpos que poseen carga eléctrica.

\end{definicion}

La fuerza eléctrica es una magnitud vectorial.

Por tanto, para especificarla completamente debemos conocer:

\[
\boxed{
\text{magnitud}
+
\text{dirección}
+
\text{sentido}
}
\]

y no únicamente un valor numérico.

Representaremos una fuerza eléctrica mediante

\[
\vec{F}.
\]

Su unidad en el Sistema Internacional es el newton:

\[
[\vec{F}]=\si{\newton}.
\]

\begin{idea}

El signo de una carga no representa la dirección de una fuerza.

La carga eléctrica es una magnitud escalar con signo, mientras que la fuerza
eléctrica es una magnitud vectorial.

La dirección y el sentido de la fuerza deben determinarse a partir de la
geometría del sistema y del carácter atractivo o repulsivo de la
interacción.

\end{idea}


% ============================================================
% MODELO DE CARGA PUNTUAL
% ============================================================

\section{Modelo de carga puntual}

La ley de Coulomb se formula de manera directa para
\textbf{cargas puntuales}.

\begin{definicion}[Carga puntual]

Una \textbf{carga puntual} es un modelo físico en el cual consideramos que
toda la carga eléctrica de un cuerpo se encuentra concentrada en un único
punto del espacio.

\end{definicion}

Ningún objeto macroscópico real tiene tamaño exactamente cero.

Por ello, la carga puntual es una idealización.

Sin embargo, esta aproximación resulta adecuada cuando las dimensiones de
los cuerpos cargados son muy pequeñas comparadas con la distancia que los
separa.

Si un cuerpo tiene una dimensión característica \(L\) y la distancia entre
los cuerpos es \(r\), la aproximación de carga puntual es especialmente
adecuada cuando

\[
L\ll r.
\]

\begin{interpretacion}

La expresión \(L\ll r\) no representa una igualdad exacta.

Indica que el tamaño del cuerpo debe ser suficientemente pequeño comparado
con la distancia relevante para que su extensión espacial pueda despreciarse
en el modelo.

\end{interpretacion}

En algunas situaciones, una distribución esféricamente simétrica de carga
también puede comportarse externamente, para el cálculo de la fuerza, como si
su carga estuviera concentrada en su centro.

Por ello, cuando trabajamos con pequeñas esferas cargadas y su radio es mucho
menor que la separación entre ellas, suele ser razonable tratarlas como
cargas puntuales.


% ============================================================
% EXPERIMENTOS DE COULOMB
% ============================================================

\section{Dependencia de la fuerza eléctrica}

Charles-Augustin de Coulomb estudió experimentalmente la interacción entre
cuerpos cargados mediante una balanza de torsión.

Sus experimentos permitieron establecer cómo depende la fuerza eléctrica de:

\begin{itemize}
    \item la magnitud de las cargas;
    \item la distancia entre ellas.
\end{itemize}


% ============================================================
% DEPENDENCIA CON LAS CARGAS
% ============================================================

\subsection{Dependencia con las cargas}

Consideremos dos cargas puntuales \(q_1\) y \(q_2\).

Experimentalmente se encuentra que, manteniendo fija la distancia entre
ellas, la magnitud de la fuerza aumenta proporcionalmente con la magnitud de
cada carga.

Así,

\[
F\propto |q_1|
\]

y también

\[
F\propto |q_2|.
\]

Combinando ambas relaciones,

\[
F\propto |q_1q_2|.
\]

Esto significa, por ejemplo, que si duplicamos únicamente una de las cargas,

\[
q_1\longrightarrow2q_1,
\]

entonces

\[
F\longrightarrow2F.
\]

Si duplicamos ambas cargas,

\[
q_1\longrightarrow2q_1,
\qquad
q_2\longrightarrow2q_2,
\]

el producto de las cargas se multiplica por cuatro:

\[
(2q_1)(2q_2)=4q_1q_2,
\]

por lo que

\[
\boxed{
F\longrightarrow4F
}.
\]


% ============================================================
% DEPENDENCIA CON LA DISTANCIA
% ============================================================

\subsection{Dependencia con la distancia}

La fuerza eléctrica también depende de la distancia \(r\) que separa las
cargas.

Experimentalmente se encuentra que

\[
F\propto\frac{1}{r^2}.
\]

Esta relación se denomina una
\textbf{dependencia inversa con el cuadrado de la distancia}.

Por ejemplo, si la distancia se duplica,

\[
r\longrightarrow2r,
\]

entonces

\[
F'
\propto
\frac{1}{(2r)^2}
=
\frac{1}{4r^2}.
\]

Por tanto,

\[
\boxed{
F'=\frac{F}{4}
}.
\]

Si la distancia se triplica,

\[
r\longrightarrow3r,
\]

entonces

\[
\boxed{
F'=\frac{F}{9}
}.
\]

Si la distancia se reduce a la mitad,

\[
r\longrightarrow\frac{r}{2},
\]

entonces

\[
F'
\propto
\frac{1}{(r/2)^2}
=
\frac{4}{r^2},
\]

por lo que

\[
\boxed{
F'=4F
}.
\]

\begin{idea}

La dependencia \(1/r^2\) es muy importante.

Duplicar la distancia no reduce la fuerza a la mitad:

\[
\boxed{
r\to2r
\quad\Longrightarrow\quad
F\to\frac{F}{4}
}.
\]

\end{idea}


% ============================================================
% LEY DE COULOMB
% ============================================================

\section{Ley de Coulomb}

Combinando la dependencia con las cargas y con la distancia,

\[
F\propto
\frac{|q_1q_2|}{r^2}.
\]

Para convertir esta proporcionalidad en una igualdad introducimos una
constante \(k\).

\begin{principio}[Ley de Coulomb]

La magnitud de la fuerza eléctrica entre dos cargas puntuales en reposo es
directamente proporcional al producto de las magnitudes de las cargas e
inversamente proporcional al cuadrado de la distancia que las separa.

\end{principio}

Matemáticamente,

\begin{equation}
    \boxed{
    F
    =
    k
    \frac{|q_1q_2|}{r^2}
    }.
    \label{eq:ley_coulomb_magnitud}
\end{equation}

donde

\begin{align*}
    F   &= \text{magnitud de la fuerza eléctrica},\\
    q_1 &= \text{primera carga eléctrica},\\
    q_2 &= \text{segunda carga eléctrica},\\
    r   &= \text{distancia entre las cargas},\\
    k   &= \text{constante de Coulomb}.
\end{align*}


% ============================================================
% VALOR ABSOLUTO
% ============================================================

\subsection{¿Por qué aparece el valor absoluto?}

En la ecuación

\[
F
=
k\frac{|q_1q_2|}{r^2},
\]

la cantidad \(F\) representa la \textbf{magnitud} de una fuerza.

Una magnitud vectorial no puede tener una magnitud negativa.

Por eso utilizamos

\[
|q_1q_2|.
\]

Por ejemplo, si

\[
q_1>0
\]

y

\[
q_2<0,
\]

el producto algebraico

\[
q_1q_2<0.
\]

Esto no significa que la magnitud de la fuerza sea negativa.

El signo del producto nos permite reconocer que las cargas tienen signos
opuestos y que, por tanto, la interacción es atractiva.

La magnitud sigue siendo

\[
F>0.
\]

\begin{errorfrecuente}

No escribir una ``magnitud negativa'' de fuerza como consecuencia de obtener

\[
q_1q_2<0.
\]

Si estamos calculando únicamente la magnitud mediante

\[
F=k\frac{|q_1q_2|}{r^2},
\]

el resultado es siempre no negativo.

El carácter atractivo o repulsivo se determina por separado.

\end{errorfrecuente}


% ============================================================
% CONSTANTE DE COULOMB
% ============================================================

\subsection{Constante de Coulomb}

Para cargas en el vacío, la constante de Coulomb tiene aproximadamente el
valor

\begin{equation}
    \boxed{
    k
    \approx
    \SI{8.99e9}{\newton\metre\squared\per\coulomb\squared}
    }.
    \label{eq:constante_coulomb}
\end{equation}

En cálculos donde no se requiere gran precisión suele utilizarse

\[
\boxed{
k
\approx
\SI{9.0e9}{\newton\metre\squared\per\coulomb\squared}
}.
\]

Las unidades de \(k\) pueden verificarse a partir de la ley de Coulomb.

Despejando,

\[
k
=
\frac{Fr^2}{|q_1q_2|}.
\]

Por tanto,

\[
[k]
=
\frac{\si{\newton}\,\si{\metre\squared}}
{\si{\coulomb\squared}},
\]

es decir,

\[
\boxed{
[k]
=
\si{\newton\metre\squared\per\coulomb\squared}
}.
\]


% ============================================================
% PERMITIVIDAD DEL VACÍO
% ============================================================

\subsection{Permitividad eléctrica del vacío}

La constante de Coulomb también puede expresarse mediante otra constante
fundamental de la electrostática denominada
\textbf{permitividad eléctrica del vacío}, representada por

\[
\varepsilon_0.
\]

La relación es

\begin{equation}
    \boxed{
    k
    =
    \frac{1}{4\pi\varepsilon_0}
    }.
    \label{eq:k_epsilon_0}
\end{equation}

El valor de la permitividad eléctrica del vacío es aproximadamente

\begin{equation}
    \boxed{
    \varepsilon_0
    \approx
    \SI{8.8541878188e-12}{\farad\per\metre}
    }.
    \label{eq:permitividad_vacio}
\end{equation}

Por tanto, la ley de Coulomb también puede escribirse como

\begin{equation}
    \boxed{
    F
    =
    \frac{1}{4\pi\varepsilon_0}
    \frac{|q_1q_2|}{r^2}
    }.
    \label{eq:coulomb_epsilon_0}
\end{equation}

Ambas expresiones son equivalentes en el vacío.

\begin{idea}

Para cálculos sencillos suele ser más práctico utilizar \(k\).

La expresión con \(\varepsilon_0\) será especialmente útil cuando las
ecuaciones electrostáticas comiencen a escribirse de una forma más general.

\end{idea}


% ============================================================
% CONDICIONES DE APLICACIÓN
% ============================================================

\subsection{Condiciones de aplicación}

La forma

\[
F
=
k\frac{|q_1q_2|}{r^2}
\]

se aplica directamente cuando:

\begin{itemize}
    \item las cargas pueden modelarse como puntuales;
    \item las cargas se encuentran en reposo dentro del modelo
    electrostático;
    \item \(r\) representa la distancia entre las dos cargas;
    \item se trabaja en el vacío, o en aire cuando la aproximación al vacío
    resulta adecuada.
\end{itemize}

Cuando existe materia entre las cargas, la interacción puede modificarse
debido a la respuesta eléctrica del material.

Por tanto, no debemos utilizar automáticamente el valor de \(k\) del vacío
para cualquier medio material sin analizar primero el modelo correspondiente.

\begin{advertencia}

La ley de Coulomb para dos cargas puntuales no debe aplicarse sin análisis a
cuerpos extensos cuya forma, tamaño o distribución de carga sean importantes.

En esos casos, el cuerpo debe modelarse adecuadamente antes de realizar el
cálculo.

\end{advertencia}


% ============================================================
% DIRECCIÓN Y SENTIDO
% ============================================================

\section{Dirección y sentido de la fuerza eléctrica}

La ecuación escalar de Coulomb proporciona la magnitud de la fuerza.

Pero una fuerza es un vector.

Debemos determinar además su dirección y su sentido.

Para dos cargas puntuales, la fuerza eléctrica actúa siempre a lo largo de
la recta que une las cargas.

Por tanto,

\[
\boxed{
\text{dirección de }\vec{F}
=
\text{recta que une }q_1\text{ y }q_2
}.
\]


% ============================================================
% CARGAS DEL MISMO SIGNO
% ============================================================

\subsection{Cargas del mismo signo}

Si

\[
q_1q_2>0,
\]

las cargas tienen el mismo signo.

La interacción es repulsiva.

Por ejemplo,

\[
+\;q_1
\qquad\qquad
+\;q_2
\]

produce fuerzas dirigidas alejando una carga de la otra:

\[
\longleftarrow
\quad
+\;q_1
\qquad\qquad
+\;q_2
\quad
\longrightarrow
\]

Lo mismo ocurre para dos cargas negativas.

Por tanto,

\[
\boxed{
q_1q_2>0
\Longrightarrow
\text{repulsión}
}.
\]


% ============================================================
% CARGAS DE SIGNO OPUESTO
% ============================================================

\subsection{Cargas de signos opuestos}

Si

\[
q_1q_2<0,
\]

las cargas tienen signos opuestos.

La interacción es atractiva.

Por ejemplo,

\[
+\;q_1
\qquad\qquad
-\;q_2
\]

produce fuerzas dirigidas una hacia la otra:

\[
+\;q_1
\quad
\longrightarrow
\qquad
\longleftarrow
\quad
-\;q_2.
\]

Por tanto,

\[
\boxed{
q_1q_2<0
\Longrightarrow
\text{atracción}
}.
\]


% ============================================================
% TERCERA LEY
% ============================================================

\section{Ley de Coulomb y tercera ley de Newton}

Supongamos que \(q_1\) ejerce una fuerza sobre \(q_2\).

La representaremos como

\[
\vec{F}_{1\to2}.
\]

Al mismo tiempo, \(q_2\) ejerce una fuerza sobre \(q_1\):

\[
\vec{F}_{2\to1}.
\]

Estas fuerzas tienen:

\begin{itemize}
    \item la misma magnitud;
    \item la misma línea de acción;
    \item sentidos opuestos.
\end{itemize}

Por tanto,

\begin{equation}
    \boxed{
    \vec{F}_{1\to2}
    =
    -\vec{F}_{2\to1}
    }.
    \label{eq:tercera_ley_coulomb}
\end{equation}

Esta relación es una manifestación de la tercera ley de Newton.

En magnitud,

\[
\boxed{
F_{1\to2}
=
F_{2\to1}
}.
\]

\begin{idea}

Aunque las cargas sean diferentes,

\[
|q_1|\neq|q_2|,
\]

las fuerzas que ejercen mutuamente tienen la misma magnitud.

Por ejemplo, una carga de

\[
\SI{1}{\micro\coulomb}
\]

y otra de

\[
\SI{100}{\micro\coulomb}
\]

ejercen fuerzas de igual magnitud entre sí.

No es correcto afirmar que la carga mayor ejerce una fuerza mayor sobre la
carga menor.

\end{idea}

\begin{errorfrecuente}

Las fuerzas

\[
\vec{F}_{1\to2}
\]

y

\[
\vec{F}_{2\to1}
\]

no se cancelan entre sí porque actúan sobre cuerpos diferentes.

Solo se suman fuerzas que actúan sobre el mismo cuerpo cuando queremos
calcular la fuerza neta sobre ese cuerpo.

\end{errorfrecuente}


% ============================================================
% EJEMPLO BÁSICO
% ============================================================

\section{Aplicación directa de la ley de Coulomb}

\begin{ejemplo}

Dos cargas puntuales

\[
q_1=+\SI{2.0}{\micro\coulomb}
\]

y

\[
q_2=-\SI{3.0}{\micro\coulomb}
\]

se encuentran separadas una distancia

\[
r=\SI{0.40}{\metre}.
\]

Determine la magnitud de la fuerza eléctrica y establezca si la interacción
es atractiva o repulsiva.

\end{ejemplo}

Primero convertimos las cargas al SI:

\[
q_1
=
+\SI{2.0e-6}{\coulomb},
\]

\[
q_2
=
-\SI{3.0e-6}{\coulomb}.
\]

La magnitud de la fuerza es

\[
F
=
k\frac{|q_1q_2|}{r^2}.
\]

Sustituyendo,

\[
F
=
\left(
\SI{8.99e9}{\newton\metre\squared\per\coulomb\squared}
\right)
\frac{
\left|
(\SI{2.0e-6}{\coulomb})
(-\SI{3.0e-6}{\coulomb})
\right|
}{
(\SI{0.40}{\metre})^2
}.
\]

Calculando,

\[
F
\approx
\SI{0.337}{\newton}.
\]

Por tanto,

\[
\boxed{
F\approx\SI{0.34}{\newton}
}.
\]

Como las cargas tienen signos opuestos,

\[
q_1q_2<0,
\]

la fuerza es de atracción.

\begin{interpretacion}

La ley de Coulomb nos dio la magnitud.

El signo de las cargas nos permitió determinar el carácter de la
interacción.

Ambas partes son necesarias para describir correctamente la fuerza.

\end{interpretacion}


% ============================================================
% RAZONAMIENTO PROPORCIONAL
% ============================================================

\section{Razonamiento proporcional con la ley de Coulomb}

No todos los problemas requieren sustituir valores numéricos.

Muchas situaciones pueden resolverse comparando estados iniciales y finales.

Partimos de

\[
F
=
k\frac{|q_1q_2|}{r^2}.
\]

Si las cargas y la distancia cambian,

\[
F'
=
k\frac{|q_1'q_2'|}{r'^2}.
\]

Dividiendo ambas expresiones,

\begin{equation}
    \boxed{
    \frac{F'}{F}
    =
    \frac{|q_1'q_2'|}{|q_1q_2|}
    \frac{r^2}{r'^2}
    }.
    \label{eq:razon_coulomb}
\end{equation}

Esta forma permite analizar rápidamente cómo cambia la fuerza.


% ============================================================
% EJEMPLO PROPORCIONAL
% ============================================================

\subsection{Ejemplo de proporcionalidad}

\begin{ejemplo}

Dos cargas ejercen inicialmente una fuerza de magnitud \(F\).

Si una carga se duplica, la otra se triplica y la distancia entre ellas se
duplica, ¿cuál será la nueva magnitud de la fuerza?

\end{ejemplo}

Tenemos

\[
q_1'=2q_1,
\]

\[
q_2'=3q_2,
\]

y

\[
r'=2r.
\]

Entonces,

\[
\frac{F'}{F}
=
\frac{|(2q_1)(3q_2)|}{|q_1q_2|}
\frac{r^2}{(2r)^2}.
\]

Simplificando,

\[
\frac{F'}{F}
=
6\left(\frac{1}{4}\right).
\]

Por tanto,

\[
\frac{F'}{F}
=
\frac{3}{2},
\]

y

\[
\boxed{
F'=\frac{3}{2}F
}.
\]

La nueva fuerza es \(1.5\) veces la original.

\begin{notapedagogica}

Este tipo de problema es útil para comprobar si se comprende realmente la
dependencia funcional

\[
F\propto\frac{|q_1q_2|}{r^2}.
\]

Conviene resolverlo primero mediante factores de cambio antes de utilizar
valores numéricos.

\end{notapedagogica}


% ============================================================
% FORMA VECTORIAL
% ============================================================

\section{Forma vectorial de la ley de Coulomb}

La expresión

\[
F
=
k\frac{|q_1q_2|}{r^2}
\]

solo proporciona la magnitud de la fuerza.

Para describir completamente la interacción necesitamos una ecuación
vectorial.

Supongamos que queremos calcular la fuerza que la carga \(q_1\) ejerce sobre
la carga \(q_2\).

Definimos un vector unitario

\[
\hat{\mathbf{r}}_{12}
\]

dirigido desde \(q_1\) hacia \(q_2\).

Entonces,

\begin{equation}
    \boxed{
    \vec{F}_{1\to2}
    =
    k
    \frac{q_1q_2}{r_{12}^{2}}
    \hat{\mathbf{r}}_{12}
    }.
    \label{eq:coulomb_vectorial}
\end{equation}

donde

\begin{align*}
    \vec{F}_{1\to2}
    &=
    \text{fuerza que }q_1\text{ ejerce sobre }q_2,
    \\
    r_{12}
    &=
    \text{distancia entre }q_1\text{ y }q_2,
    \\
    \hat{\mathbf{r}}_{12}
    &=
    \text{vector unitario dirigido de }q_1\text{ hacia }q_2.
\end{align*}


% ============================================================
% SIGNO EN FORMA VECTORIAL
% ============================================================

\subsection{Significado del signo en la forma vectorial}

Existe una diferencia fundamental entre

\[
F
=
k\frac{|q_1q_2|}{r^2}
\]

y

\[
\vec{F}_{1\to2}
=
k\frac{q_1q_2}{r^2}
\hat{\mathbf{r}}_{12}.
\]

En la ecuación de la \textbf{magnitud} utilizamos

\[
|q_1q_2|.
\]

En la ecuación \textbf{vectorial} conservamos el signo de

\[
q_1q_2.
\]

Si

\[
q_1q_2>0,
\]

entonces

\[
\vec{F}_{1\to2}
\]

tiene el mismo sentido que

\[
\hat{\mathbf{r}}_{12},
\]

lo cual corresponde a repulsión.

Si

\[
q_1q_2<0,
\]

el factor escalar es negativo, por lo que el vector fuerza apunta en sentido
opuesto a

\[
\hat{\mathbf{r}}_{12},
\]

lo cual corresponde a atracción.

\begin{idea}

En la forma vectorial, el signo no significa que exista una ``fuerza
negativa''.

Indica la orientación del vector respecto del vector unitario elegido.

\end{idea}


% ============================================================
% POSICIONES VECTORIALES
% ============================================================

\subsection{Forma vectorial a partir de las posiciones}

Supongamos que las cargas se encuentran en las posiciones

\[
\vec{r}_1
\]

y

\[
\vec{r}_2.
\]

El vector que va desde \(q_1\) hasta \(q_2\) es

\begin{equation}
    \boxed{
    \vec{r}_{12}
    =
    \vec{r}_2-\vec{r}_1
    }.
    \label{eq:vector_posicion_12}
\end{equation}

Su magnitud es

\[
r_{12}
=
\abs{\vec{r}_{12}},
\]

y el vector unitario correspondiente es

\begin{equation}
    \boxed{
    \hat{\mathbf{r}}_{12}
    =
    \frac{\vec{r}_{12}}{r_{12}}
    }.
    \label{eq:unitario_12}
\end{equation}

Por tanto,

\[
\vec{F}_{1\to2}
=
k
\frac{q_1q_2}{r_{12}^{2}}
\frac{\vec{r}_{12}}{r_{12}}.
\]

Así obtenemos una forma especialmente útil:

\begin{equation}
    \boxed{
    \vec{F}_{1\to2}
    =
    kq_1q_2
    \frac{
    \vec{r}_2-\vec{r}_1
    }{
    \abs{\vec{r}_2-\vec{r}_1}^{3}
    }
    }.
    \label{eq:coulomb_posiciones}
\end{equation}

Esta expresión incorpora simultáneamente:

\begin{itemize}
    \item la magnitud;
    \item la dirección;
    \item el sentido;
    \item las posiciones de las cargas.
\end{itemize}


% ============================================================
% SUPERPOSICIÓN
% ============================================================

\section{Principio de superposición}

La ley de Coulomb describe directamente la interacción entre dos cargas
puntuales.

Sin embargo, muchos sistemas contienen tres, cuatro o muchas más cargas.

Cuando varias cargas actúan simultáneamente sobre una carga determinada,
cada una produce su propia fuerza.

Experimentalmente se encuentra que estas fuerzas se combinan mediante suma
vectorial.

\begin{principio}[Principio de superposición]

La fuerza eléctrica neta sobre una carga es la suma vectorial de las fuerzas
eléctricas individuales ejercidas sobre ella por todas las demás cargas.

\end{principio}

Si queremos determinar la fuerza sobre una carga \(q_i\),

\begin{equation}
    \boxed{
    \vec{F}_i
    =
    \sum_{j\neq i}
    \vec{F}_{j\to i}
    }.
    \label{eq:superposicion_coulomb}
\end{equation}

Por ejemplo, si tres cargas \(q_1\), \(q_2\) y \(q_3\) ejercen fuerzas sobre
una carga \(q_0\),

\begin{equation}
    \boxed{
    \vec{F}_{\text{neta}}
    =
    \vec{F}_{1\to0}
    +
    \vec{F}_{2\to0}
    +
    \vec{F}_{3\to0}
    }.
\end{equation}

Cada una de estas fuerzas se calcula como si las otras cargas no estuvieran
presentes.

Después se realiza la suma vectorial.


% ============================================================
% MAGNITUDES NO SE SUMAN SIEMPRE
% ============================================================

\subsection{La superposición es vectorial}

Este punto es fundamental.

En general,

\[
\boxed{
F_{\text{neta}}
\neq
F_1+F_2+\cdots
}
\]

si \(F_1,F_2,\ldots\) representan únicamente magnitudes.

Lo que debemos sumar son los vectores:

\[
\boxed{
\vec{F}_{\text{neta}}
=
\vec{F}_1+\vec{F}_2+\cdots
}.
\]

Solo cuando todas las fuerzas se encuentran sobre una misma recta podemos
realizar una suma algebraica utilizando signos previamente definidos para
representar sentidos opuestos.


% ============================================================
% CASO UNIDIMENSIONAL
% ============================================================

\section{Superposición en una dimensión}

Cuando todas las cargas están ubicadas sobre una misma línea, el problema
vectorial puede reducirse a una suma algebraica.

Primero se elige un eje y un sentido positivo.

Por ejemplo,

\[
+x
\quad\text{hacia la derecha}.
\]

Entonces:

\[
F_x>0
\]

representará una fuerza hacia la derecha y

\[
F_x<0
\]

una fuerza hacia la izquierda.

\begin{advertencia}

El signo de una componente de fuerza depende del sistema de coordenadas
elegido.

No depende directamente del signo de la carga.

Una carga negativa puede experimentar una fuerza hacia \(+x\) o hacia
\(-x\), dependiendo de la configuración de las demás cargas.

\end{advertencia}


% ============================================================
% EJEMPLO 1D
% ============================================================

\subsection{Ejemplo: tres cargas alineadas}

\begin{ejemplo}

Una carga

\[
q_0=+\SI{1.0}{\micro\coulomb}
\]

se encuentra en \(x=0\).

Una carga

\[
q_1=+\SI{4.0}{\micro\coulomb}
\]

se encuentra en

\[
x=-\SI{0.30}{\metre},
\]

y una carga

\[
q_2=-\SI{9.0}{\micro\coulomb}
\]

se encuentra en

\[
x=+\SI{0.60}{\metre}.
\]

Determine la fuerza eléctrica neta sobre \(q_0\).

\end{ejemplo}

Tomamos como positivo el sentido hacia la derecha.

Primero analizamos la fuerza producida por \(q_1\).

Como \(q_1\) y \(q_0\) son positivas, se repelen.

Como \(q_1\) está a la izquierda de \(q_0\), la repulsión empuja a \(q_0\)
hacia la derecha.

Por tanto,

\[
F_{1\to0,x}>0.
\]

Su magnitud es

\[
F_{1\to0}
=
k
\frac{|q_1q_0|}
{r_{10}^{2}}.
\]

Sustituyendo,

\[
F_{1\to0}
=
\left(8.99\times10^9\right)
\frac{
(4.0\times10^{-6})
(1.0\times10^{-6})
}{
(0.30)^2
}.
\]

Entonces,

\[
F_{1\to0}
\approx
\SI{0.400}{\newton}.
\]

Así,

\[
F_{1\to0,x}
=
+\SI{0.400}{\newton}.
\]

Ahora analizamos \(q_2\).

Como \(q_2\) es negativa y \(q_0\) positiva, se atraen.

Como \(q_2\) está a la derecha, la fuerza sobre \(q_0\) apunta hacia la
derecha.

Entonces,

\[
F_{2\to0,x}>0.
\]

Su magnitud es

\[
F_{2\to0}
=
k
\frac{|q_2q_0|}
{r_{20}^{2}}.
\]

Sustituyendo,

\[
F_{2\to0}
=
\left(8.99\times10^9\right)
\frac{
(9.0\times10^{-6})
(1.0\times10^{-6})
}{
(0.60)^2
}.
\]

De donde

\[
F_{2\to0}
\approx
\SI{0.225}{\newton}.
\]

Entonces,

\[
F_{2\to0,x}
=
+\SI{0.225}{\newton}.
\]

Aplicamos superposición:

\[
F_{\text{neta},x}
=
F_{1\to0,x}
+
F_{2\to0,x}.
\]

Por tanto,

\[
F_{\text{neta},x}
=
\SI{0.400}{\newton}
+
\SI{0.225}{\newton}.
\]

Finalmente,

\[
\boxed{
F_{\text{neta},x}
\approx
+\SI{0.625}{\newton}
}.
\]

El signo positivo indica que la fuerza neta apunta hacia la derecha.

\begin{interpretacion}

Aunque una interacción es repulsiva y la otra atractiva, ambas fuerzas
apuntan en este caso hacia la derecha.

Por eso no debemos decidir los signos de las componentes únicamente mirando
si la interacción es atractiva o repulsiva.

Primero debemos dibujar o imaginar la dirección real de cada fuerza.

\end{interpretacion}


% ============================================================
% SUPERPOSICIÓN EN 2D
% ============================================================

\section{Superposición en dos dimensiones}

Cuando las fuerzas no se encuentran sobre una misma línea, debemos realizar
una suma vectorial.

Una forma sistemática consiste en descomponer cada fuerza en componentes
cartesianas.

Si una fuerza \(\vec{F}\) forma un ángulo \(\theta\) medido desde el eje
positivo \(x\),

\[
F_x
=
F\cos\theta,
\]

\[
F_y
=
F\sin\theta.
\]

La fuerza neta tiene componentes

\begin{equation}
    \boxed{
    F_{\text{neta},x}
    =
    \sum_i F_{i,x}
    }
\end{equation}

y

\begin{equation}
    \boxed{
    F_{\text{neta},y}
    =
    \sum_i F_{i,y}
    }.
\end{equation}

Por tanto,

\begin{equation}
    \boxed{
    \vec{F}_{\text{neta}}
    =
    F_{\text{neta},x}\,\hat{\mathbf{i}}
    +
    F_{\text{neta},y}\,\hat{\mathbf{j}}
    }.
\end{equation}

Su magnitud es

\begin{equation}
    \boxed{
    F_{\text{neta}}
    =
    \sqrt{
    F_{\text{neta},x}^{\,2}
    +
    F_{\text{neta},y}^{\,2}
    }
    }.
    \label{eq:magnitud_fuerza_neta}
\end{equation}


% ============================================================
% DIRECCIÓN RESULTANTE
% ============================================================

\subsection{Dirección de la fuerza resultante}

Si conocemos las componentes,

\[
F_x
\]

y

\[
F_y,
\]

podemos obtener un ángulo de referencia mediante

\[
\tan\theta
=
\frac{|F_y|}{|F_x|}.
\]

Sin embargo, el ángulo final debe ubicarse en el cuadrante correcto de
acuerdo con los signos de

\[
F_x
\]

y

\[
F_y.
\]

\begin{advertencia}

La expresión

\[
\theta
=
\tan^{-1}
\left(
\frac{F_y}{F_x}
\right)
\]

puede producir un ángulo que no identifica correctamente el cuadrante si se
utiliza sin analizar los signos de las componentes.

Primero determine físicamente el cuadrante del vector resultante.

\end{advertencia}


% ============================================================
% EJEMPLO 2D
% ============================================================

\subsection{Ejemplo: suma de fuerzas perpendiculares}

\begin{ejemplo}

Una carga \(q=+\SI{2.0}{\micro\coulomb}\) se encuentra en el origen.
Una carga \(q_1=+\SI{3.0}{\micro\coulomb}\) se encuentra en
\((\SI{0.40}{\metre},0)\), y otra carga
\(q_2=+\SI{4.0}{\micro\coulomb}\) se encuentra en
\((0,\SI{0.30}{\metre})\).

Determine la fuerza neta sobre la carga situada en el origen.

\end{ejemplo}

Como las tres cargas son positivas, \(q_1\) repele a \(q\) hacia la izquierda
y \(q_2\) la repele hacia abajo.

\begin{center}
\begin{tikzpicture}[>=Latex,scale=1.1]

\draw[->] (-2.6,0) -- (5.1,0) node[right] {$x$};
\draw[->] (0,-3.8) -- (0,4.1) node[above] {$y$};

\draw (4,0.10) -- (4,-0.10) node[below=4pt] {\(\SI{0.40}{\metre}\)};
\draw (0.10,3) -- (-0.10,3) node[left=4pt] {\(\SI{0.30}{\metre}\)};

\filldraw (0,0) circle (2.2pt) node[above left=3pt] {$q$};
\filldraw (4,0) circle (2.2pt) node[above=5pt] {$q_1$};
\filldraw (0,3) circle (2.2pt) node[right=5pt] {$q_2$};

\draw[->,very thick] (0,0) -- (-1.3,0)
node[midway,above=4pt] {$F_1q$};

\draw[->,very thick] (0,0) -- (0,-3.0)
node[midway,right=5pt] {$F_2q$};

\draw[->,thick,dashed] (0,0) -- (-1.3,-3.0)
node[pos=0.78,left=4pt] {$F_{\text{neta}}$};

\draw[densely dotted] (-1.3,-3.0) -- (-1.3,0);
\draw[densely dotted] (-1.3,-3.0) -- (0,-3.0);

\end{tikzpicture}
\end{center}

La magnitud de la fuerza producida por \(q_1\) es

\[
F_1=k\frac{|q_1q|}{r_1^2}.
\]

Sustituyendo,

\[
F_1=(8.99\times10^9)
\frac{(3.0\times10^{-6})(2.0\times10^{-6})}{(0.40)^2}
\approx\SI{0.337}{\newton}.
\]

Como esta fuerza es horizontal y apunta hacia la izquierda,

\[
\boxed{F_{1x}=-\SI{0.337}{\newton}}, \qquad \boxed{F_{1y}=0}.
\]

Para \(q_2\),

\[
F_2=k\frac{|q_2q|}{r_2^2}.
\]

Sustituyendo,

\[
F_2=(8.99\times10^9)
\frac{(4.0\times10^{-6})(2.0\times10^{-6})}{(0.30)^2}
\approx\SI{0.799}{\newton}.
\]

Como esta fuerza es vertical y apunta hacia abajo,

\[
\boxed{F_{2x}=0}, \qquad \boxed{F_{2y}=-\SI{0.799}{\newton}}.
\]

Aplicando superposición en cada eje,

\[
F_x=F_{1x}+F_{2x}=-0.337+0,
\qquad
F_y=F_{1y}+F_{2y}=0-0.799.
\]

Por tanto,

\[
\boxed{F_x=-\SI{0.337}{\newton}},
\qquad
\boxed{F_y=-\SI{0.799}{\newton}}.
\]

La magnitud de la fuerza neta es

\[
F_{\text{neta}}=\sqrt{F_x^2+F_y^2}.
\]

Sustituyendo,

\[
F_{\text{neta}}
=\sqrt{(-0.337)^2+(-0.799)^2}
\approx\SI{0.867}{\newton}.
\]

Por tanto,

\[
\boxed{F_{\text{neta}}\approx\SI{0.87}{\newton}}.
\]

Como \(F_x<0\) y \(F_y<0\), la fuerza resultante se encuentra en el tercer
cuadrante. El ángulo respecto del eje \(-x\) es

\[
\theta=\tan^{-1}\left(\frac{|F_y|}{|F_x|}\right)
=\tan^{-1}\left(\frac{0.799}{0.337}\right)
\approx67.1^\circ.
\]

Así,

\[
\boxed{\theta\approx67.1^\circ\text{ por debajo del eje }-x}.
\]

En conclusión,

\[
\boxed{
F_x=-\SI{0.337}{\newton},\qquad
F_y=-\SI{0.799}{\newton},\qquad
F_{\text{neta}}\approx\SI{0.87}{\newton}
}.
\]


% ============================================================
% SIGNOS DE LAS COMPONENTES
% ============================================================

\section{Signo de las componentes de una fuerza}

Una fuente frecuente de errores consiste en confundir:

\begin{itemize}
    \item el signo de una carga;
    \item el signo del producto \(q_1q_2\);
    \item el signo de una componente de fuerza.
\end{itemize}

Son conceptos diferentes.

Por ejemplo, si una fuerza apunta hacia la izquierda y elegimos

\[
+x
\]

hacia la derecha, entonces

\[
F_x<0.
\]

Esto ocurre independientemente de que las cargas involucradas sean positivas
o negativas.

De igual manera, si una fuerza forma un ángulo \(\theta\) con el eje positivo
\(x\),

\[
F_x=F\cos\theta.
\]

Si el vector se encuentra en una región donde

\[
\cos\theta<0,
\]

la componente horizontal será negativa.

\begin{idea}

El signo de una componente proviene de la orientación del vector respecto del
sistema de coordenadas.

No debe colocarse un signo negativo adicional solamente porque alguna de las
cargas sea negativa.

\end{idea}


% ============================================================
% MÉTODO GENERAL
% ============================================================

\section{Método para resolver problemas con varias cargas}

Una estrategia ordenada reduce considerablemente los errores en problemas de
ley de Coulomb.

\subsection{1. Identificar la carga sobre la que se pide la fuerza}

Primero determine con claridad cuál es la carga de interés.

Todas las fuerzas que sumaremos deben actuar sobre esa misma carga.

\subsection{2. Dibujar la configuración}

Ubique las cargas, sus signos y las distancias relevantes.

Si existen coordenadas, establezca los ejes antes de calcular.

\subsection{3. Convertir al Sistema Internacional}

En particular,

\[
\SI{1}{\micro\coulomb}
=
\SI{1e-6}{\coulomb},
\]

\[
\SI{1}{\nano\coulomb}
=
\SI{1e-9}{\coulomb}.
\]

Las distancias deben expresarse en metros cuando utilizamos

\[
k
=
\SI{8.99e9}{\newton\metre\squared\per\coulomb\squared}.
\]

\subsection{4. Analizar cada interacción por separado}

Antes de calcular componentes, determine si cada interacción es atractiva o
repulsiva.

Dibuje mental o explícitamente cada vector de fuerza.

\subsection{5. Calcular la magnitud de cada fuerza}

Para cada par de cargas,

\[
F_i
=
k
\frac{|q_iq|}{r_i^2}.
\]

\subsection{6. Descomponer en componentes}

Cuando sea necesario,

\[
F_{i,x}
=
F_i\cos\theta_i,
\]

\[
F_{i,y}
=
F_i\sin\theta_i,
\]

asignando los signos de acuerdo con el sentido real del vector.

\subsection{7. Aplicar superposición}

Sume las componentes:

\[
F_{\text{neta},x}
=
\sum_i F_{i,x},
\]

\[
F_{\text{neta},y}
=
\sum_i F_{i,y}.
\]

\subsection{8. Calcular magnitud y dirección}

Finalmente,

\[
F_{\text{neta}}
=
\sqrt{
F_{\text{neta},x}^{\,2}
+
F_{\text{neta},y}^{\,2}
}.
\]

Determine la dirección utilizando las componentes y verificando el cuadrante.

\subsection{9. Interpretar el resultado}

Antes de aceptar la respuesta, compruebe:

\begin{itemize}
    \item que las unidades sean newtons;
    \item que la dirección obtenida sea físicamente razonable;
    \item que las fuerzas atractivas y repulsivas estén orientadas
    correctamente;
    \item que la magnitud tenga un orden de magnitud razonable.
\end{itemize}


% ============================================================
% SIMETRÍA
% ============================================================

\section{Uso de la simetría}

Algunas configuraciones permiten simplificar considerablemente la suma de
fuerzas.

Supongamos que dos cargas iguales se encuentran simétricamente ubicadas
respecto de una tercera carga.

Si ambas ejercen fuerzas de igual magnitud,

\[
F_1=F_2,
\]

puede ocurrir que ciertas componentes se cancelen.

Por ejemplo,

\[
F_{1,y}
=
-F_{2,y},
\]

por lo que

\[
F_{\text{neta},y}=0.
\]

Mientras tanto, las componentes horizontales pueden sumarse:

\[
F_{\text{neta},x}
=
F_{1,x}+F_{2,x}.
\]

\begin{idea}

Antes de realizar cálculos extensos, observe siempre la geometría del
sistema.

La simetría puede mostrar inmediatamente qué componentes se cancelan y cuáles
se refuerzan.

\end{idea}


% ============================================================
% EQUILIBRIO
% ============================================================

\section{Equilibrio eléctrico de una carga}

Una carga puede encontrarse bajo la acción de varias fuerzas eléctricas.

Si la suma vectorial de todas ellas es cero,

\[
\boxed{
\sum_i\vec{F}_i=0
},
\]

la carga tiene fuerza eléctrica neta nula.

Esto no significa necesariamente que no existan fuerzas.

Puede ocurrir que varias fuerzas no nulas se cancelen vectorialmente.

En una dimensión, la condición puede escribirse como

\[
\boxed{
\sum_iF_{i,x}=0
}.
\]

\begin{interpretacion}

Una fuerza neta igual a cero no significa ausencia de interacción.

Significa que las diferentes fuerzas que actúan sobre el cuerpo producen una
resultante nula.

\end{interpretacion}


% ============================================================
% COMPARACIÓN CON GRAVITACIÓN
% ============================================================

\section{Comparación cualitativa con la interacción gravitatoria}

La ley de Coulomb presenta una estructura matemática semejante a la ley de
gravitación universal:

\[
F_{\text{eléctrica}}
\propto
\frac{|q_1q_2|}{r^2},
\]

mientras que

\[
F_{\text{gravitatoria}}
\propto
\frac{m_1m_2}{r^2}.
\]

En ambos casos la magnitud disminuye con el cuadrado de la distancia.

Sin embargo, existe una diferencia fundamental.

La interacción gravitatoria entre masas ordinarias es atractiva.

La interacción eléctrica puede ser:

\[
\boxed{\text{atractiva}}
\]

o

\[
\boxed{\text{repulsiva}},
\]

porque existen dos signos de carga eléctrica.

Además, a escala microscópica las interacciones eléctricas pueden ser
enormemente más intensas que las gravitatorias.


% ============================================================
% ERRORES FRECUENTES
% ============================================================

\section{Errores y confusiones frecuentes}

\begin{errorfrecuente}

\textbf{Usar microcoulombs o nanocoulombs directamente con la constante
\(k\).}

Si se utiliza

\[
k=
\SI{8.99e9}{\newton\metre\squared\per\coulomb\squared},
\]

las cargas deben expresarse en coulombs.

\end{errorfrecuente}


\begin{errorfrecuente}

\textbf{Olvidar elevar la distancia al cuadrado.}

La ley de Coulomb contiene

\[
\frac{1}{r^2},
\]

no

\[
\frac{1}{r}.
\]

\end{errorfrecuente}


\begin{errorfrecuente}

\textbf{Obtener una magnitud de fuerza negativa.}

Para calcular la magnitud se utiliza

\[
F
=
k\frac{|q_1q_2|}{r^2}.
\]

El carácter atractivo o repulsivo se analiza aparte.

\end{errorfrecuente}


\begin{errorfrecuente}

\textbf{Sumar directamente las magnitudes de fuerzas que tienen direcciones
diferentes.}

El principio de superposición exige una suma vectorial:

\[
\vec{F}_{\text{neta}}
=
\sum_i\vec{F}_i.
\]

\end{errorfrecuente}


\begin{errorfrecuente}

\textbf{Asignar el signo de una componente según el signo de la carga.}

El signo de \(F_x\) o \(F_y\) depende del sentido del vector respecto de los
ejes elegidos.

\end{errorfrecuente}


\begin{errorfrecuente}

\textbf{Pensar que la carga de mayor magnitud ejerce una fuerza mayor sobre
la carga pequeña.}

Por la tercera ley de Newton,

\[
\vec{F}_{1\to2}
=
-\vec{F}_{2\to1}.
\]

Las fuerzas tienen igual magnitud y sentidos opuestos.

\end{errorfrecuente}


\begin{errorfrecuente}

\textbf{Sumar las dos fuerzas de un par de acción y reacción para obtener
cero.}

Estas dos fuerzas actúan sobre cuerpos diferentes.

No forman parte de la misma suma de fuerzas sobre un solo cuerpo.

\end{errorfrecuente}


\begin{errorfrecuente}

\textbf{Confundir distancia entre cargas con una coordenada.}

En general,

\[
r\neq x_2-x_1
\]

si las cargas no se encuentran sobre el mismo eje.

En dos o tres dimensiones debe calcularse la distancia geométrica entre las
posiciones.

Por ejemplo, en dos dimensiones,

\[
r
=
\sqrt{
(x_2-x_1)^2
+
(y_2-y_1)^2
}.
\]

\end{errorfrecuente}


\begin{errorfrecuente}

\textbf{Aplicar automáticamente la ley de Coulomb para cargas puntuales a
cualquier cuerpo extenso.}

Debe comprobarse primero si el cuerpo puede aproximarse razonablemente por
una carga puntual o si su distribución espacial de carga debe considerarse.

\end{errorfrecuente}


% ============================================================
% PREGUNTAS PARA COMPROBAR COMPRENSIÓN
% ============================================================

\section{Preguntas para comprobar comprensión}

\begin{enumerate}

    \item Si se duplica una de las cargas y todo lo demás permanece igual,
    ¿cómo cambia la fuerza eléctrica?

    \item Si se duplican ambas cargas, ¿cómo cambia la fuerza?

    \item Si la distancia entre dos cargas se duplica, ¿cómo cambia la
    fuerza?

    \item Si la distancia se reduce a la mitad, ¿cómo cambia la fuerza?

    \item ¿Por qué se utiliza \(\abs{q_1q_2}\) cuando calculamos la magnitud
    de la fuerza?

    \item ¿Qué información aporta el signo de \(q_1q_2\)?

    \item ¿En qué dirección actúa la fuerza eléctrica entre dos cargas
    puntuales?

    \item ¿Las fuerzas que dos cargas ejercen mutuamente tienen siempre la
    misma magnitud?

    \item Si una carga es cien veces mayor que otra, ¿ejerce una fuerza cien
    veces mayor sobre la carga pequeña?

    \item ¿Por qué las fuerzas de acción y reacción no se cancelan al
    estudiar el movimiento de una sola carga?

    \item ¿Qué diferencia existe entre sumar magnitudes y sumar vectores?

    \item ¿Cuándo podemos tratar una suma de fuerzas eléctricas como una
    suma algebraica?

    \item ¿Qué principio utilizamos cuando varias cargas actúan
    simultáneamente sobre una misma carga?

    \item ¿De qué depende el signo de una componente \(F_x\)?

    \item ¿Puede una carga positiva experimentar una fuerza hacia la
    izquierda?

    \item ¿Puede una carga negativa experimentar una fuerza hacia la
    derecha?

    \item ¿Qué significa físicamente que la fuerza eléctrica neta sobre una
    carga sea cero?

    \item ¿Por qué es útil analizar la simetría antes de comenzar un
    cálculo?

\end{enumerate}


% ============================================================
% RESUMEN
% ============================================================

\begin{resumen}

\begin{itemize}

    \item La fuerza eléctrica es una magnitud vectorial que describe la
    interacción entre cuerpos cargados.

    \item Para dos cargas puntuales, la magnitud de la fuerza viene dada por

    \[
    \boxed{
    F
    =
    k\frac{|q_1q_2|}{r^2}
    }.
    \]

    \item En el vacío,

    \[
    \boxed{
    k
    \approx
    \SI{8.99e9}{\newton\metre\squared\per\coulomb\squared}
    }.
    \]

    \item La constante de Coulomb se relaciona con la permitividad del vacío
    mediante

    \[
    \boxed{
    k
    =
    \frac{1}{4\pi\varepsilon_0}
    }.
    \]

    \item La magnitud de la fuerza es directamente proporcional al producto
    de las magnitudes de las cargas:

    \[
    F\propto|q_1q_2|.
    \]

    \item La fuerza es inversamente proporcional al cuadrado de la distancia:

    \[
    F\propto\frac{1}{r^2}.
    \]

    \item Si las cargas tienen el mismo signo, la interacción es repulsiva.

    \item Si las cargas tienen signos opuestos, la interacción es atractiva.

    \item La fuerza actúa a lo largo de la recta que une las dos cargas.

    \item Las fuerzas que dos cargas ejercen mutuamente cumplen

    \[
    \boxed{
    \vec{F}_{1\to2}
    =
    -\vec{F}_{2\to1}
    }.
    \]

    \item Una forma vectorial de la ley de Coulomb es

    \[
    \boxed{
    \vec{F}_{1\to2}
    =
    k
    \frac{q_1q_2}{r_{12}^{2}}
    \hat{\mathbf{r}}_{12}
    }.
    \]

    \item Cuando varias cargas actúan sobre una misma carga se aplica el
    principio de superposición:

    \[
    \boxed{
    \vec{F}_i
    =
    \sum_{j\neq i}
    \vec{F}_{j\to i}
    }.
    \]

    \item En dos dimensiones, la superposición puede realizarse mediante
    componentes:

    \[
    F_{\text{neta},x}
    =
    \sum_iF_{i,x},
    \]

    \[
    F_{\text{neta},y}
    =
    \sum_iF_{i,y}.
    \]

    \item La magnitud de la fuerza resultante es

    \[
    \boxed{
    F_{\text{neta}}
    =
    \sqrt{
    F_{\text{neta},x}^{\,2}
    +
    F_{\text{neta},y}^{\,2}
    }
    }.
    \]

    \item Los signos de las componentes dependen del sistema de coordenadas
    y de la orientación de los vectores, no directamente del signo de las
    cargas.

\end{itemize}

\end{resumen}